\documentclass[12pt,a4paper]{article}
\usepackage{ugmath}
\newcommand{\paren}[1]{\left( #1 \right)}
\newcommand{\pd}[3]{\left(\frac{\partial #1}{\partial #2}\right)_{#3}}
\newcommand{\alumno}{Ricardo León Martínez}
\newcommand{\materia}{Termodinámica}
\newcommand{\tarea}{Tarea 2}
\newcommand{\fecha}{16/02/2024}

\begin{document}
    \begin{center}
    {\large\textbf{UNIVERSIDAD DE GUANAJUATO}}\\[0.3cm]
    {\normalsize\textbf{DIVISIÓN DE CIENCIAS E INGENIERÍA FÍSICA}}\\
    {\normalsize\textbf{CAMPUS GUANAJUATO}}\\[1cm]

    {\Large\textbf{\tarea\ (\materia)}}\\[1cm]
    \end{center}

    \textbf{Nombre:} \alumno \hfill
    \textbf{Profesor:} Dr. Néstor M. De Los Santos López \hfill
    \textbf{Fecha:} \fecha \\[0.3cm]

    \begin{problema}[Diferencial total]
        Investigar si $df$ representa un diferencial total:
        \begin{enumerate}
            \item[(1)] $df=\cos x\sin y\,dx-\sin x\cos y\,dy$.
            \item[(2)] $df=\sin x\cos y\,dx+\cos x\sin y\,dy$.
            \item[(3)] $df=x^{3}y^{2}dx-y^{3}x^{2}dy$.
        \end{enumerate}
    \end{problema}

    \begin{solucion}
        Una expresión $df=M(x,y)dx+N(x,y)dy$, con $M,N$ de clase $C^{1}$, es un diferencial
        total (exacto) si y sólo si
        \begin{equation*}
            \frac{\partial M}{\partial y}=\frac{\partial N}{\partial x}.
        \end{equation*}
        \begin{enumerate}
            \item[(1)] Aquí $M=\cos x\sin y$, $N=-\sin x\cos y$. Entonces
            \begin{equation*}
                \frac{\partial M}{\partial y}=\cos x\cos y, \qquad \frac{\partial N}{\partial x}=-\cos x\cos y.
            \end{equation*}
            Como $\cos x\cos y\neq-\cos x\cos y$ en general, \textbf{$df$ no es un diferencial total.}
            \item[(2)] Aquí $M=\sin x\cos y$, $N=\cos x\sin y$. Entonces
            \begin{equation*}
                \frac{\partial M}{\partial y}=-\sin x\sin y=\frac{\partial N}{\partial x}.
            \end{equation*}
            Se cumple la igualdad, así que \textbf{$df$ sí es un diferencial total.} De hecho
            \begin{equation*}
                f(x,y)=-\cos x\cos y+\text{cte.}
            \end{equation*}
            satisface $df=\sin x\cos y\,dx+\cos x\sin y\,dy$.
            \item[(3)] Aquí $M=x^{3}y^{2}$, $N=-x^{2}y^{3}$. Entonces
            \begin{equation*}
                \frac{\partial M}{\partial y}=2x^{3}y, \qquad \frac{\partial N}{\partial x}=-2xy^{3}.
            \end{equation*}
            Como $2x^{3}y\neq-2xy^{3}$ en general, \textbf{$df$ no es un diferencial total.}
        \end{enumerate}
    \end{solucion}

    \begin{problema}[Derivadas parciales]
        Sea $f(x,y,z)=0$. Verifique las siguientes relaciones:
        \begin{enumerate}
            \item[(1)] $\displaystyle\pd{x}{y}{z}=\frac{1}{\pd{y}{x}{z}}$.
            \item[(2)] $\displaystyle\pd{x}{y}{z}\pd{y}{z}{x}\pd{z}{x}{y}=-1$.
        \end{enumerate}
    \end{problema}

    \begin{solucion}
        La relación $f(x,y,z)=0$ define implícitamente a cada variable en función de las
        otras dos. Escribamos, por brevedad, $f_{x}\equiv\pd{f}{x}{y,z}$, $f_{y}\equiv\pd{f}{y}{x,z}$,
        $f_{z}\equiv\pd{f}{z}{x,y}$, y supongamos $f_{x},f_{y},f_{z}\neq0$ (condición del teorema
        de la función implícita). El diferencial de $f=0$ es
        \begin{equation*}
            f_{x}\,dx+f_{y}\,dy+f_{z}\,dz=0.
        \end{equation*}
        \begin{enumerate}
            \item[(1)] Fijando $z$ (i.e., $dz=0$), la relación anterior da $f_{x}dx+f_{y}dy=0$,
            de donde
            \begin{equation*}
                \pd{x}{y}{z}=-\frac{f_{y}}{f_{x}}, \qquad \pd{y}{x}{z}=-\frac{f_{x}}{f_{y}}.
            \end{equation*}
            Multiplicando,
            \begin{equation*}
                \pd{x}{y}{z}\cdot\pd{y}{x}{z}=\paren{-\frac{f_{y}}{f_{x}}}\paren{-\frac{f_{x}}{f_{y}}}=1,
            \end{equation*}
            lo cual equivale exactamente a la relación pedida.
            \item[(2)] De manera análoga, fijando cada variable a su turno,
            \begin{equation*}
                \pd{x}{y}{z}=-\frac{f_{y}}{f_{x}}, \qquad
                \pd{y}{z}{x}=-\frac{f_{z}}{f_{y}}, \qquad
                \pd{z}{x}{y}=-\frac{f_{x}}{f_{z}}.
            \end{equation*}
            Multiplicando las tres,
            \begin{equation*}
                \pd{x}{y}{z}\pd{y}{z}{x}\pd{z}{x}{y}
                =(-1)^{3}\cdot\frac{f_{y}}{f_{x}}\cdot\frac{f_{z}}{f_{y}}\cdot\frac{f_{x}}{f_{z}}
                =-1.
            \end{equation*}
        \end{enumerate}
    \end{solucion}

    \begin{problema}[Integral de camino]
        \begin{enumerate}
            \item[a)] La integral de camino
            \begin{equation*}
                I(C)=\int_{A(C)}^{B}\left[\alpha(x,y)dx+\beta(x,y)dy\right]
            \end{equation*}
            debe calcularse para dos puntos finales fijos $A$ y $B$ en el plano $xy$. En
            general, tendrá diferentes valores para diferentes caminos $C_{i}$ entre $A$ y $B$.
            Demostrar que $I$ es independiente del camino si y sólo si $\dfrac{\partial\alpha}{\partial y}=\dfrac{\partial\beta}{\partial x}$.
            \item[b)] Elija especialmente $\alpha(x,y)=y^{2}e^{x}$, $\beta(x,y)=2ye^{x}$; calcule
            \begin{equation*}
                I_{AB}=\int_{A}^{B}\left[\alpha(x,y)dx+\beta(x,y)dy\right],
            \end{equation*}
            para $A=(0,0)$, $B=(1,1)$. Considere primero si el problema es razonable. En caso
            afirmativo, encuentre el valor de $I_{AB}$.
            \item[c)] Haga las mismas investigaciones que en b), pero con los roles
            intercambiados de $\alpha$ y $\beta$, i.e., $\alpha(x,y)=2ye^{x}$, $\beta(x,y)=y^{2}e^{x}$.
        \end{enumerate}
    \end{problema}

    \begin{solucion}
        \begin{enumerate}
            \item[a)] Sean $C_{1}$ y $C_{2}$ dos caminos de $A$ a $B$ contenidos en un dominio
            $D\subset\mathbb{R}^{2}$ simplemente conexo. Supongamos primero que
            $\alpha_{y}=\beta_{x}$ en todo $D$. El camino cerrado $\Gamma=C_{1}-C_{2}$ (ir por
            $C_{1}$ de $A$ a $B$ y regresar por $C_{2}$) encierra una región $R\subset D$, y por
            el teorema de Green,
            \begin{equation*}
                \oint_{\Gamma}\left(\alpha\,dx+\beta\,dy\right)=\iint_{R}\paren{\frac{\partial\beta}{\partial x}-\frac{\partial\alpha}{\partial y}}dA=0.
            \end{equation*}
            Pero $\oint_{\Gamma}=\int_{C_{1}}-\int_{C_{2}}$, así que $\int_{C_{1}}=\int_{C_{2}}$:
            $I$ no depende del camino.

            Recíprocamente, si $I$ es independiente del camino, entonces $\oint_{\Gamma}=0$ para
            \emph{todo} lazo cerrado $\Gamma$ en $D$ (formado como diferencia de dos caminos
            entre los mismos puntos). Por Green,
            \begin{equation*}
                \iint_{R}\paren{\beta_{x}-\alpha_{y}}dA=0
            \end{equation*}
            para toda región $R\subset D$. Tomando $R$ un disco de radio $\varepsilon\to0$
            alrededor de un punto arbitrario $(x_{0},y_{0})$ y usando la continuidad de
            $\beta_{x}-\alpha_{y}$, el teorema del valor medio para integrales fuerza
            $\beta_{x}(x_{0},y_{0})-\alpha_{y}(x_{0},y_{0})=0$. Como $(x_{0},y_{0})$ es
            arbitrario, $\alpha_{y}=\beta_{x}$ en todo $D$.
            \item[b)] Verificamos primero la exactitud:
            \begin{equation*}
                \frac{\partial\alpha}{\partial y}=2ye^{x}=\frac{\partial\beta}{\partial x}.
            \end{equation*}
            Se cumple, así que por el inciso a) el problema \emph{sí es razonable}: $I_{AB}$
            no depende del camino elegido, sólo de los extremos. Buscamos $\varphi$ con
            $\varphi_{x}=y^{2}e^{x}$, $\varphi_{y}=2ye^{x}$. Integrando la primera respecto a $x$,
            \begin{equation*}
                \varphi(x,y)=y^{2}e^{x}+g(y).
            \end{equation*}
            Derivando respecto a $y$ e igualando a $\beta$, $2ye^{x}+g'(y)=2ye^{x}\implies g'(y)=0$.
            Así, $\varphi(x,y)=y^{2}e^{x}$ (salvo constante), y
            \begin{equation*}
                I_{AB}=\varphi(1,1)-\varphi(0,0)=1^{2}e^{1}-0=e.
            \end{equation*}
            \item[c)] Ahora $\alpha=2ye^{x}$, $\beta=y^{2}e^{x}$. Verificamos exactitud:
            \begin{equation*}
                \frac{\partial\alpha}{\partial y}=2e^{x}, \qquad \frac{\partial\beta}{\partial x}=y^{2}e^{x}.
            \end{equation*}
            Como $2e^{x}\neq y^{2}e^{x}$ en general, la forma $\alpha\,dx+\beta\,dy$
            \textbf{no es exacta}: por el inciso a), el problema \emph{no es razonable} tal como
            está planteado, pues $I_{AB}$ depende del camino elegido, no sólo de $A$ y $B$.
            Ilustramos esto con dos caminos concretos de $A=(0,0)$ a $B=(1,1)$:
            \begin{itemize}
                \item \emph{Recta $y=x$}, con $x=y=t$, $t\in[0,1]$:
                \begin{equation*}
                    I=\int_{0}^{1}\left[2te^{t}+t^{2}e^{t}\right]dt=\int_{0}^{1}\frac{d}{dt}\paren{t^{2}e^{t}}dt=t^{2}e^{t}\Big|_{0}^{1}=e.
                \end{equation*}
                \item \emph{Camino quebrado} $(0,0)\to(1,0)\to(1,1)$: en el primer tramo $y=0$,
                de modo que $\alpha\,dx+\beta\,dy=0$; en el segundo tramo $x=1$ fijo, $dx=0$, y
                \begin{equation*}
                    I=\int_{0}^{1}y^{2}e\,dy=\frac{e}{3}.
                \end{equation*}
            \end{itemize}
            Como $e\neq e/3$, se confirma que la integral depende del camino: $I_{AB}$ no está
            bien definida sin especificar la trayectoria.
        \end{enumerate}
    \end{solucion}

    \begin{problema}[Expansión, compresibilidad]
        Suponiendo que la ecuación de estado de un gas es $p=p(V,T)$, expresar el coeficiente
        de expansión de volumen térmico isobárico,
        \begin{equation*}
            \beta=\frac{1}{V}\pd{V}{T}{P},
        \end{equation*}
        y la compresibilidad isotérmica,
        \begin{equation*}
            \kappa_{T}=-\frac{1}{V}\pd{V}{p}{T},
        \end{equation*}
        en términos de derivadas de $p(V,T)$.
    \end{problema}

    \begin{solucion}
        Aplicamos las identidades del problema 2 a la relación implícita
        $F(p,V,T)\equiv p-p(V,T)=0$. Por la regla de reciprocidad,
        \begin{equation*}
            \pd{V}{p}{T}=\frac{1}{\pd{p}{V}{T}}.
        \end{equation*}
        Así,
        \begin{equation*}
            \boxed{\kappa_{T}=-\frac{1}{V\pd{p}{V}{T}}.}
        \end{equation*}
        Por la regla cíclica aplicada a $(V,T,p)$,
        \begin{equation*}
            \pd{V}{T}{p}\pd{T}{p}{V}\pd{p}{V}{T}=-1
            \implies
            \pd{V}{T}{p}=-\frac{1}{\pd{T}{p}{V}\pd{p}{V}{T}}
            =-\frac{\pd{p}{T}{V}}{\pd{p}{V}{T}},
        \end{equation*}
        usando de nuevo la reciprocidad $\pd{T}{p}{V}=1/\pd{p}{T}{V}$. Por lo tanto,
        \begin{equation*}
            \boxed{\beta=-\frac{1}{V}\cdot\frac{\pd{p}{T}{V}}{\pd{p}{V}{T}}.}
        \end{equation*}
        En particular, dividiendo ambas cajas, $\beta/\kappa_{T}=\pd{p}{T}{V}$, una relación de
        consistencia útil entre las tres cantidades.
    \end{solucion}

    \begin{problema}[Ecuación de estado]
        Para una sustancia homogénea con número molar $n$ se han encontrado las siguientes
        relaciones:
        \begin{equation*}
            \beta=\frac{nR}{pV}, \qquad \kappa_{T}=\frac{1}{T}+\frac{a}{V},
        \end{equation*}
        con $\beta,\kappa_{T}$ definidas como en el problema anterior y $R$ la constante
        universal de los gases. Encuentra la ecuación de estado $f(T,p,V)=0$.
    \end{problema}

    \begin{solucion}
        De $\beta=\frac1V\pd VTp=\dfrac{nR}{pV}$ obtenemos, sin depender de $\kappa_{T}$,
        \begin{equation*}
            \pd{V}{T}{p}=\frac{nR}{p}.
        \end{equation*}
        Integrando en $T$ a $p$ fijo, $V(T,p)=\dfrac{nRT}{p}+\varphi(p)$, con $\varphi$ una
        función por determinar. Ahora imponemos $\kappa_{T}=-\frac1V\pd Vpp=\frac1T+\frac aV$
        literalmente:
        \begin{equation*}
            \pd{V}{p}{T}=-\frac{V}{T}-a
            \implies
            -\frac{nRT}{p^{2}}+\varphi'(p)=-\frac{nRT/p+\varphi(p)}{T}-a
            =-\frac{nR}{p}-\frac{\varphi(p)}{T}-a.
        \end{equation*}
        Agrupando por potencias de $T$,
        \begin{equation*}
            \underbrace{\paren{-\frac{nR}{p^{2}}}}_{\text{coef.\ de }T}T+\underbrace{\paren{\varphi'(p)+\frac{nR}{p}+a}}_{\text{coef.\ constante}}+\underbrace{\varphi(p)}_{\text{coef.\ de }1/T}\cdot\frac1T=0.
        \end{equation*}
        Como $\{T,1,1/T\}$ son funciones linealmente independientes de $T$, esta identidad para
        \emph{todo} $T$ exige que el coeficiente de $T$ se anule: $-nR/p^{2}=0$, lo cual es
        imposible para $n,R,p\neq0$. Es decir, \textbf{tal como está escrito, con el término
        $1/T$, el par $(\beta,\kappa_{T})$ no es integrable}: no existe ninguna $V(T,p)$ que
        satisfaga ambas relaciones simultáneamente (falla la igualdad de derivadas cruzadas
        $\partial_{p}\partial_{T}V=\partial_{T}\partial_{p}V$).

        La forma correcta, integrable, es la que reemplaza $1/T$ por $1/p$:
        $\kappa_{T}=\dfrac1p+\dfrac aV$, dimensionalmente compatible con $\beta=nR/(pV)$, que
        también exhibe un factor $1/p$. Verifiquemos que esta versión sí es integrable: con
        $\pd Vpp=-\frac{V}{p}-a$,
        \begin{equation*}
            \partial_{p}\pd{V}{T}{p}=\partial_{p}\paren{\frac{nR}{p}}=-\frac{nR}{p^{2}},
            \qquad
            \partial_{T}\pd{V}{p}{T}=\partial_{T}\paren{-\frac Vp-a}=-\frac1p\pd{V}{T}{p}=-\frac{nR}{p^{2}}.
        \end{equation*}
        Las derivadas cruzadas coinciden idénticamente: el sistema es integrable. Procedemos con
        $\kappa_{T}=1/p+a/V$. Sustituyendo $V=\frac{nRT}{p}+\varphi(p)$,
        \begin{equation*}
            -\frac{nRT}{p^{2}}+\varphi'(p)=-\frac{nRT/p+\varphi(p)}{p}-a
            =-\frac{nRT}{p^{2}}-\frac{\varphi(p)}{p}-a.
        \end{equation*}
        Los términos $-nRT/p^{2}$ se cancelan en ambos lados, quedando la EDO lineal
        \begin{equation*}
            \varphi'(p)+\frac{\varphi(p)}{p}=-a.
        \end{equation*}
        Con factor integrante $p$, $(p\varphi)'=-ap$, de donde $p\varphi(p)=-\dfrac{ap^{2}}{2}+C$,
        es decir $\varphi(p)=-\dfrac{ap}{2}+\dfrac{C}{p}$. Tomando la constante de integración
        $C=0$ (para que, en el límite $a\to0$, se recupere exactamente la ley del gas ideal
        $V=nRT/p$), obtenemos
        \begin{equation*}
            V(T,p)=\frac{nRT}{p}-\frac{ap}{2},
        \end{equation*}
        es decir, la ecuación de estado buscada es
        \begin{equation*}
            \boxed{f(T,p,V)=pV-nRT+\frac{a}{2}p^{2}=0.}
        \end{equation*}
        (Puede verificarse por sustitución directa que esta $V(T,p)$ reproduce exactamente
        $\beta=nR/(pV)$ y $\kappa_{T}=1/p+a/V$.)
    \end{solucion}

    \begin{problema}[Ecuación de Van der Waals]
        La ecuación de Van der Waals, para una sustancia con volumen molar $V$,
        \begin{equation*}
            \paren{p+\frac{a}{V^{2}}}(V-b)=RT,
        \end{equation*}
        describe cualitativamente la transición gas-líquido.
        \begin{enumerate}
            \item[a)] Expresar las constantes $a$ y $b$ mediante $V_{c}$ y $T_{c}$.
            \item[b)] Reescribir la ecuación en términos de las cantidades reducidas
            $\pi=p/p_{c}$, $v=V/V_{c}$, $t=T/T_{c}$.
            \item[c)] Calcular $\kappa_{T}=-\frac1V\pd VpT$ para $V=V_{c}$. ¿Qué comportamiento
            muestra $\kappa_{T}$ cuando $T\to T_{c}^{+}$? ¿Cómo se entiende físicamente?
            \item[d)] Calcular $\beta=\frac1V\pd VTp$.
        \end{enumerate}
    \end{problema}

    \begin{solucion}
        Despejamos $p=\dfrac{RT}{V-b}-\dfrac{a}{V^{2}}$.
        \begin{enumerate}
            \item[a)] El punto crítico es un punto de inflexión horizontal de la isoterma
            $T=T_{c}$: $\pd pVT=0$ y $\left(\dfrac{\partial^{2}p}{\partial V^{2}}\right)_{T}=0$ (en notación abreviada, $p_{V}=0=p_{VV}$)
            en $V=V_{c}$. Calculamos
            \begin{equation*}
                p_{V}=-\frac{RT}{(V-b)^{2}}+\frac{2a}{V^{3}}, \qquad
                p_{VV}=\frac{2RT}{(V-b)^{3}}-\frac{6a}{V^{4}}.
            \end{equation*}
            Igualando ambas a cero en $(V_{c},T_{c})$ y dividiendo la segunda entre la primera,
            \begin{equation*}
                \frac{2}{V_{c}-b}=\frac{3}{V_{c}} \implies b=\frac{V_{c}}{3}.
            \end{equation*}
            Sustituyendo en $p_{V}=0$ (con $V_{c}-b=2V_{c}/3$),
            \begin{equation*}
                \frac{RT_{c}}{(2V_{c}/3)^{2}}=\frac{2a}{V_{c}^{3}} \implies a=\frac{9RT_{c}V_{c}}{8}.
            \end{equation*}
            Como bono, sustituyendo en la ecuación de estado se obtiene la razón universal
            \begin{equation*}
                p_{c}=\frac{3RT_{c}}{8V_{c}}, \qquad \text{es decir} \qquad \frac{p_{c}V_{c}}{RT_{c}}=\frac38.
            \end{equation*}
            \item[b)] Sustituyendo $p=\pi p_{c}$, $V=vV_{c}$, $T=tT_{c}$, $b=V_{c}/3$,
            $a=9RT_{c}V_{c}/8$ y $p_{c}=3RT_{c}/(8V_{c})$ en $p=RT/(V-b)-a/V^{2}$:
            \begin{align*}
                \pi p_{c}&=\frac{RtT_{c}}{V_{c}(v-1/3)}-\frac{9RT_{c}V_{c}/8}{v^{2}V_{c}^{2}}
                =\frac{RT_{c}}{V_{c}}\left[\frac{t}{v-1/3}-\frac{9}{8v^{2}}\right].
            \end{align*}
            Como $RT_{c}/V_{c}=8p_{c}/3$, dividiendo entre $p_{c}$:
            \begin{equation*}
                \pi=\frac{8}{3}\left[\frac{t}{v-1/3}-\frac{9}{8v^{2}}\right]=\frac{8t}{3v-1}-\frac{3}{v^{2}},
            \end{equation*}
            es decir, la \textbf{ecuación reducida de Van der Waals} (ley de estados
            correspondientes):
            \begin{equation*}
                \boxed{\paren{\pi+\frac{3}{v^{2}}}(3v-1)=8t.}
            \end{equation*}
            \item[c)] Usando $\kappa_{T}=-1/(Vp_{V})$ del problema 4, evaluamos $p_{V}$ en
            $V=V_{c}$ (a temperatura $T$ general, no necesariamente $T_{c}$), con
            $V_{c}-b=2V_{c}/3$ y $a=9RT_{c}V_{c}/8$:
            \begin{align*}
                p_{V}\big|_{V_{c}}&=-\frac{RT}{(2V_{c}/3)^{2}}+\frac{2a}{V_{c}^{3}}
                =-\frac{9RT}{4V_{c}^{2}}+\frac{9RT_{c}}{4V_{c}^{2}}
                =\frac{9R}{4V_{c}^{2}}(T_{c}-T).
            \end{align*}
            Por lo tanto,
            \begin{equation*}
                \boxed{\kappa_{T}(V_{c})=-\frac{1}{V_{c}\cdot\frac{9R}{4V_{c}^{2}}(T_{c}-T)}=\frac{4V_{c}}{9R(T-T_{c})}.}
            \end{equation*}
            Cuando $T\to T_{c}^{+}$, $\kappa_{T}\to+\infty$: la compresibilidad isotérmica
            \emph{diverge} en el punto crítico. Físicamente, esto significa que cerca de $T_{c}$
            una variación arbitrariamente pequeña de presión produce una variación de volumen
            arbitrariamente grande; las fluctuaciones de densidad del fluido se vuelven muy
            grandes, lo cual es precisamente el origen microscópico de la opalescencia crítica.
            \item[d)] Usando $\beta=-\frac1V\cdot p_{T}/p_{V}$ del problema 4, con
            $p_{T}=\pd pTV=\dfrac{R}{V-b}$ y $p_{V}$ como arriba (ahora en $V$ general),
            \begin{align*}
                \beta&=-\frac1V\cdot\frac{R/(V-b)}{-RT/(V-b)^{2}+2a/V^{3}}
                =\frac{RV^{2}(V-b)}{RTV^{3}-2a(V-b)^{2}}.
            \end{align*}
            Es decir,
            \begin{equation*}
                \boxed{\beta=\frac{RV^{2}(V-b)}{RTV^{3}-2a(V-b)^{2}}.}
            \end{equation*}
            Como verificación de consistencia, en el límite diluido $V\to\infty$ el término
            dominante da $\beta\to RV^{3}/(RTV^{3})=1/T$, el valor del gas ideal, como debe ser.
        \end{enumerate}
    \end{solucion}

    \begin{problema}[Gas real (Dieterici)]
        La ecuación de estado de un gas real viene dada por
        \begin{equation*}
            p=\frac{nRT}{V-nb}\,e^{-na/(RTV)},
        \end{equation*}
        con $n=N/N_{A}$, $R$ la constante universal de los gases y $a,b$ constantes materiales
        (gas de Dieterici).
        \begin{enumerate}
            \item[1)] A partir de la expansión virial en la densidad de partículas $\rho=N/V$,
            \begin{equation*}
                p=k_{B}T\rho\paren{1+\sum_{v=1}^{\infty}B_{v}\rho^{v}},
            \end{equation*}
            calcular $B_{1}$ y la temperatura de Boyle $T_{B}$ (definida por $B_{1}(T_{B})=0$),
            en términos de $a$ y $b$.
            \item[2)] Comparar $B_{1}$ con el coeficiente virial de la ecuación de Van der Waals.
            ¿Qué significado físico tienen $a$ y $b$?
            \item[3)] ¿Cómo se relaciona $\pd p\rho T$ con $\kappa_{T}$? ¿Qué signo debe
            esperarse por argumentos físicos?
            \item[4)] Calcular $\pd p\rho T$ para el gas de Dieterici, hallar $T_{0}(\rho)$
            donde este cociente se anula, y $T_{c}$ como el máximo de $T_{0}(\rho)$. Expresar
            $a,b$ mediante $T_{c},\rho_{c}$. ¿Qué relación existe entre $T_{c}$ y $T_{B}$?
            \item[5)] Describir cualitativamente las isotermas $p(\rho)$. ¿Dónde no son físicas?
            ¿Cuándo se licúa el gas al aumentar la presión?
        \end{enumerate}
    \end{problema}

    \begin{solucion}
        \begin{enumerate}
            \item[1)] Con $V=N/\rho$ y $N=nN_{A}$, la ecuación de estado se reescribe en
            $(\rho,T)$: usando $V-nb=n(N_{A}/\rho-b)$,
            \begin{equation*}
                p=\frac{RT}{N_{A}/\rho-b}\,e^{-a\rho/(RTN_{A})}
                =k_{B}T\cdot\frac{\rho}{1-b\rho/N_{A}}\cdot e^{-a\rho/(RTN_{A})},
            \end{equation*}
            usando $R/N_{A}=k_{B}$; nótese que esta expresión ya no depende de $n$, como debe
            ser para una cantidad intensiva. Para aligerar la notación definamos las constantes
            por partícula
            \begin{equation*}
                \gamma\equiv\frac{b}{N_{A}}, \qquad \alpha\equiv\frac{a}{N_{A}^{2}},
            \end{equation*}
            de modo que $a\rho/(RTN_{A})=\alpha\rho/(k_{B}T)$ y
            \begin{equation*}
                p=k_{B}T\rho\cdot\frac{1}{1-\gamma\rho}\cdot e^{-\alpha\rho/(k_{B}T)}.
            \end{equation*}
            Expandiendo a primer orden en $\rho$,
            \begin{equation*}
                \frac{p}{k_{B}T\rho}=\left[1+\gamma\rho+O(\rho^{2})\right]\left[1-\frac{\alpha}{k_{B}T}\rho+O(\rho^{2})\right]
                =1+\paren{\gamma-\frac{\alpha}{k_{B}T}}\rho+O(\rho^{2}),
            \end{equation*}
            de donde
            \begin{equation*}
                \boxed{B_{1}=\gamma-\frac{\alpha}{k_{B}T}=\frac{1}{N_{A}}\paren{b-\frac{a}{RT}}.}
            \end{equation*}
            $B_{1}(T_{B})=0$ da
            \begin{equation*}
                \boxed{T_{B}=\frac{a}{Rb}.}
            \end{equation*}
            \item[2)] Para Van der Waals,
            \begin{equation*}
                p=\frac{RT}{V-b}-\frac{a}{V^{2}}
                =\frac{RT}{V}\paren{1+\frac{b}{V}+\cdots}-\frac{a}{V^{2}}
                =\frac{RT}{V}\left[1+\paren{b-\frac{a}{RT}}\frac1V+\cdots\right],
            \end{equation*}
            de modo que su primer coeficiente virial (en volumen molar) es también
            $b-a/(RT)$: \textbf{ambas ecuaciones de estado coinciden hasta el segundo
            coeficiente virial}, y sólo difieren en coeficientes de orden superior. En ambos
            modelos, $b$ representa el volumen excluido por el tamaño finito de las moléculas
            (repulsión de corto alcance) y $a$ mide la intensidad de la atracción intermolecular
            de largo alcance (reduce la presión respecto al gas ideal).
            \item[3)] Como $\rho=N/V$ (a $N$ fijo), $V=N/\rho$ y $\pd V\rho{N,T}=-N/\rho^{2}=-V/\rho$.
            Por la regla de la cadena,
            \begin{equation*}
                \pd p\rho T=\pd pVT\cdot\pd V\rho T=\paren{-\frac{1}{V\kappa_{T}}}\paren{-\frac{V}{\rho}}=\frac{1}{\rho\kappa_{T}},
            \end{equation*}
            usando $\pd pVT=-1/(V\kappa_{T})$ (problema 4). Es decir,
            \begin{equation*}
                \boxed{\kappa_{T}=\frac{1}{\rho\,\pd p\rho T}.}
            \end{equation*}
            La estabilidad mecánica de la materia exige $\kappa_{T}>0$ (si no, el sistema
            colapsaría o explotaría espontáneamente); como $\rho>0$, debe cumplirse
            $\pd p\rho T>0$: comprimir el gas a temperatura fija (aumentar $\rho$) debe aumentar
            la presión.
            \item[4)] De $\ln p=\ln(k_{B}T)+\ln\rho-\ln(1-\gamma\rho)-\alpha\rho/(k_{B}T)$,
            \begin{equation*}
                \pd p\rho T=p\left[\frac1\rho+\frac{\gamma}{1-\gamma\rho}-\frac{\alpha}{k_{B}T}\right].
            \end{equation*}
            Como $p>0$, esto se anula cuando el corchete se anula, lo que define $T_{0}(\rho)$:
            \begin{equation*}
                \frac1\rho+\frac{\gamma}{1-\gamma\rho}=\frac{\alpha}{k_{B}T_{0}}
                \implies
                k_{B}T_{0}=\frac{\alpha\rho(1-\gamma\rho)}{(1-\gamma\rho)+\gamma\rho}=\alpha\rho(1-\gamma\rho).
            \end{equation*}
            Así,
            \begin{equation*}
                T_{0}(\rho)=\frac{\alpha}{k_{B}}\rho(1-\gamma\rho),
            \end{equation*}
            una parábola invertida en $\rho$, nula en $\rho=0$ y en $\rho=1/\gamma$. Su máximo
            se obtiene de $dT_{0}/d\rho=0$: $1-2\gamma\rho=0\implies\rho_{c}=1/(2\gamma)$, y
            \begin{equation*}
                T_{c}=T_{0}(\rho_{c})=\frac{\alpha}{4k_{B}\gamma}.
            \end{equation*}
            Invirtiendo estas relaciones,
            \begin{equation*}
                \gamma=\frac{1}{2\rho_{c}} \implies \boxed{b=\frac{N_{A}}{2\rho_{c}},}
                \qquad
                \alpha=4k_{B}\gamma T_{c}=\frac{2k_{B}T_{c}}{\rho_{c}} \implies \boxed{a=\frac{2N_{A}RT_{c}}{\rho_{c}}.}
            \end{equation*}
            Comparando con $T_{B}=\alpha/(k_{B}\gamma)$ y $T_{c}=\alpha/(4k_{B}\gamma)$,
            \begin{equation*}
                \boxed{T_{B}=4T_{c}.}
            \end{equation*}
            \item[5)] Para $T>T_{c}$, $\pd p\rho T>0$ en todo $\rho$ físico: la isoterma
            $p(\rho)$ es monótona creciente (comportamiento de fluido supercrítico, una sola
            fase). Para $T<T_{c}$, $T_{0}(\rho)$ tiene dos raíces reales y $p(\rho)$ desarrolla
            un máximo y un mínimo locales (un "lazo" tipo van der Waals): en el tramo entre
            ambos, $\pd p\rho T<0$, lo cual viola $\kappa_{T}>0$ y por tanto \textbf{no es
            físico} (mecánicamente inestable); ese tramo se sustituye, en la práctica, por una
            línea horizontal de coexistencia líquido-gas (construcción de Maxwell) a la presión
            de saturación correspondiente. En $T=T_{c}$ el máximo y el mínimo se funden en un
            único punto de inflexión horizontal: el punto crítico. Por lo tanto, \textbf{sólo
            para $T<T_{c}$} se puede licuar el gas aumentando la presión (cruzando la región de
            coexistencia); para $T>T_{c}$ no existe una transición de fase discontinua al
            comprimir.
        \end{enumerate}
    \end{solucion}

    \begin{problema}[Gas ideal vs. gas real]
        Las desviaciones del comportamiento de un gas ideal se aproximan mediante:
        \begin{enumerate}
            \item[1)] Volumen excluido: $p(V-nb)=nRT$.
            \item[2)] Virial en presión: $pV=nRT(1+A_{1}p)$, $A_{1}=A_{1}(T)$.
            \item[3)] Virial en volumen: $pV=nRT\paren{1+\dfrac{B_{1}}{V}}$, $B_{1}=B_{1}(T)$.
        \end{enumerate}
        Calcule, para las tres, $\kappa_{T}$ y $\beta$, y compare con el gas ideal.
    \end{problema}

    \begin{solucion}
        Para el gas ideal, $V=nRT/p$, de donde $\beta_{\text{id}}=1/T$ y $\kappa_{T,\text{id}}=1/p$.
        \begin{enumerate}
            \item[1)] De $V=nb+\dfrac{nRT}{p}$,
            \begin{equation*}
                \pd VTp=\frac{nR}{p}, \qquad \pd Vpp=-\frac{nRT}{p^{2}}.
            \end{equation*}
            Como $V=n(RT+bp)/p$,
            \begin{equation*}
                \boxed{\beta=\frac{nR/p}{V}=\frac{R}{RT+bp},}\qquad
                \boxed{\kappa_{T}=\frac{nRT/p^{2}}{V}=\frac{RT}{p(RT+bp)}.}
            \end{equation*}
            Ambas son menores que sus análogas ideales ($1/T$, $1/p$): el volumen excluido hace
            al gas menos expansible y menos compresible que el ideal.
            \item[2)] Aquí $V=\dfrac{nRT}{p}+nRTA_{1}(T)$. A $p$ fijo,
            \begin{equation*}
                \pd VTp=\frac{nR}{p}+nR\left[A_{1}+TA_{1}'\right],
            \end{equation*}
            y dividiendo entre $V=\frac{nRT}p(1+pA_1)$,
            \begin{equation*}
                \boxed{\beta=\frac1T+\frac{pA_{1}'(T)}{1+pA_{1}(T)}}
            \end{equation*}
            (si $A_{1}$ no depende de $T$, se recupera $\beta=1/T$ exactamente). A $T$ fijo,
            $A_{1}(T)$ es constante y
            \begin{equation*}
                \pd Vpp=-\frac{nRT}{p^{2}},
            \end{equation*}
            (el término $nRTA_{1}$ no depende de $p$), de donde
            \begin{equation*}
                \boxed{\kappa_{T}=\frac{nRT/p^{2}}{V}=\frac{1}{p\left(1+pA_{1}(T)\right)}.}
            \end{equation*}
            \item[3)] Aquí conviene usar $p=p(V,T)$ del problema 4, ya que $V$ no se despeja
            fácilmente:
            \begin{equation*}
                p=\frac{nRT}{V}+\frac{nRTB_{1}(T)}{V^{2}}.
            \end{equation*}
            Entonces
            \begin{align*}
                p_{T}&=\frac{nR}{V}+\frac{nR\left[B_{1}+TB_{1}'\right]}{V^{2}}, &
                p_{V}&=-\frac{nRT}{V^{2}}\left(1+\frac{2B_{1}}{V}\right).
            \end{align*}
            Usando $\kappa_{T}=-1/(Vp_{V})$ y $\beta=-p_{T}/(Vp_{V})$ del problema 4,
            \begin{equation*}
                \boxed{\kappa_{T}=\frac1p\cdot\frac{1+B_{1}/V}{1+2B_{1}/V},}\qquad
                \boxed{\beta=\frac1T\cdot\frac{1+\dfrac{B_{1}+TB_{1}'}{V}}{1+\dfrac{2B_{1}}{V}}.}
            \end{equation*}
            A orden dominante en $1/V$ (gas diluido), $\kappa_{T}\approx\dfrac1p\paren{1-\dfrac{B_{1}}{V}}$
            y $\beta\approx\dfrac1T\paren{1+\dfrac{TB_{1}'-B_{1}}{TV}}$: en los tres modelos, las
            correcciones al comportamiento ideal ($1/p$, $1/T$) son de orden $1/V$, como se
            espera de una expansión virial en la densidad.
        \end{enumerate}
    \end{solucion}

    \begin{problema}[Un capacitor]
        Las placas de un capacitor de placas paralelas con capacidad $C$ llevan las cargas $Q$
        y $-Q$. Cuando las placas se acercan un poco, la capacidad aumenta un $dC$.
        \begin{enumerate}
            \item[a)] ¿Qué trabajo mecánico debe realizarse con ese pequeño acercamiento?
            \item[b)] ¿Cómo cambia la energía del campo en el capacitor?
            \item[c)] ¿Es reversible el cambio de estado?
        \end{enumerate}
    \end{problema}

    \begin{solucion}
        El capacitor está aislado (carga $Q$ fija, sin conexión a una fuente externa), así que
        no hay trabajo eléctrico de transporte de carga ($\delta W_{\text{elec}}=V\,dq=0$ pues
        $dq=0$); el único intercambio de energía es el trabajo mecánico $\delta W$ realizado al
        mover las placas. La energía electrostática almacenada es
        \begin{equation*}
            U=\frac{Q^{2}}{2C}.
        \end{equation*}
        \begin{enumerate}
            \item[a)-b)] Con $Q$ fijo,
            \begin{equation*}
                dU=\frac{Q^{2}}{2}\,d\paren{\frac1C}=-\frac{Q^{2}}{2C^{2}}\,dC.
            \end{equation*}
            Por conservación de energía (sin intercambio de calor, proceso puramente mecánico
            cuasiestático), $\delta W=dU$, es decir
            \begin{equation*}
                \boxed{\delta W=-\frac{Q^{2}}{2C^{2}}\,dC.}
            \end{equation*}
            Como el enunciado da $dC>0$ (las placas se acercan), $\delta W<0$: el agente
            externo \emph{no} necesita suministrar trabajo positivo; por el contrario, las
            placas cargadas se atraen entre sí, y para que el acercamiento sea cuasiestático
            (no una caída libre) el agente externo debe \emph{extraer} (absorber) una cantidad
            de trabajo $|\delta W|=Q^{2}dC/(2C^{2})$, que las placas realizan sobre él. La
            energía del campo, en consecuencia, \textbf{disminuye} en exactamente esa cantidad,
            $dU=-Q^{2}dC/(2C^{2})<0$: la energía perdida por el campo es la energía extraída
            mecánicamente.
            \item[c)] Sí, el proceso es \textbf{reversible}: al tratarse de un acercamiento
            cuasiestático sin disipación (sin corrientes resistivas ni fricción), basta con
            invertir el movimiento —separar las placas lentamente— para recorrer exactamente el
            mismo camino en sentido inverso, devolviendo al campo la energía $Q^{2}dC/(2C^{2})$
            mediante el trabajo mecánico ahora suministrado por el agente externo, y regresando
            al sistema a su estado original sin generación neta de entropía.
        \end{enumerate}
    \end{solucion}

\end{document}
